\section{Методы коллаборативной фильтрации}

В данном разделе рассмотрены методы рекомендаций рецептов, учитывающие только оценки пользователей.
Также приведён метод гибридной фильтрации, учитывающий ещё и содержание рецептов.

\subsection{Простая коллаборативная фильтрация}
Коллаборативная фильтрация позволяет предсказать оценку, которую даст выбранный пользователь какому-либо рецепту, на основе оценок других пользователей~\cite{segaran}.
Оценка пользователя представляет собой число в пределах $[a, b]$.

Пусть:
\begin{itemize}
    \item $N$ -- количество пользователей;
    \item $R_i$ -- множество рецептов, оценённых пользователем $u_i$, $i \in \overline{1, N}$;
    \item $u'$ -- пользователь, оценки которого необходимо предсказывать;
    \item $R'$ -- множество рецептов, которое оценил пользователь $u'$.
\end{itemize}

Вычисляется коэффициент подобия между пользователями $u'$ и $u_i$ для $i \in \overline{1, N}$.
Для этого выполняются следующие действия:
\begin{itemize}
    \item находится пересечение $R$ множеств $R'$ и $R_i$;
    \item собираются векторы оценок $a, b$, где $a_k, b_k$ -- оценки пользователей $u', u$ соответственно для рецепта $R_k, k \in \overline{1, |R|}$;
\end{itemize}
Таким образом, $\gamma = F(a, b)$.
В качестве метрики $F$ обычно выступает Евклидово расстояние или коэффициент корреляции Пирсона~\cite{segaran}.
Евклидово расстояние вычисляется по формуле~\cite{segaran}:
\begin{equation}
    F(a, b) = \sqrt{\sum\limits_{i=1}^{n} (a_i - b_i)^2},
    \label{eq:dist}
\end{equation}
а коэффициент корреляции Пирсона -- по формуле~\cite{segaran}:
\begin{equation}
    F(a, b) = \dfrac{\sum\limits_{i=1}^{n} a_i b_i - \dfrac{\left(\sum\limits_{i=1}^{n} a_i\right) \left(\sum\limits_{i=1}^{n} b_i\right)}{n}}{\sqrt{\left(\sum\limits_{i=1}^{n} a_i^2 - \dfrac{\left(\sum\limits_{i=1}^{n} a_i\right)^2}{n}\right) \left(\sum\limits_{i=1}^{n} b_i^2 - \dfrac{\left(\sum\limits_{i=1}^{n} b_i\right)^2}{n}\right)}}
\end{equation}
где $n = |a| = |b|$.

Пусть:
\begin{itemize}
    \item $\alpha_k, k \in \overline{1, M}$ -- оценка пользователя $u'$ для каждого рецепта, где $M$ -- количество рецептов;
    \item $I_k$ -- множество номеров пользователей, оценивших рецепт $k, k \in \overline{1, M}$;
    \item $\beta_{k,i}$ -- оценка рецепта $k$ пользователем $i$, $k \in \overline{1, M}, i \in I_k$;
    \item $\alpha_k$ -- оценка пользователя $u'$ для рецепта c номером $k$.
\end{itemize}
Тогда $\alpha_k$ рассчитывается по формуле~\cite{segaran}:
\begin{equation}
    \alpha_k = \frac{\sum\limits_{i \in I_k} (\gamma_i\beta_{k,i})}{\sum\limits_{i \in I_k} (\gamma_i)}.
\end{equation}
Таким образом на основе оценок других пользователей предсказывается оценка выбранного пользователя.

К особенностям данного метода относится необходимость в хранении матрицы оценок рецептов, где элемент в строке $i$ и столбце $j$ содержит либо оценку пользователя $i$ для рецепта $j$, либо метку о том, что оценка отсутствует.
Заполнение такой матрицы является трудоёмкой задачей и должно происходить заранее~\cite{segaran}.

Коллаборативная фильтрация имеет следующий ряд особенностей~[1]:
\begin{itemize}
    \item нельзя применить к пользователю, данные о котором отсутствуют;
    \item используется только матрица предпочтений и данные об оценках пользователя;
    \item метод не позволяет выявить новых предпочтений у пользователя, так как основан на истории его оценок.
\end{itemize}

\subsection{Матричные разложения}
В рассматриваемом методе делается предположение, что пользователи и объекты описываются вектором скрытых признаков, представленных в виде чисел.
Причём скалярное произведение вектора пользователя на вектор объекта даёт оценку, которую данный пользователь поставил этому объекту~\cite{yandex_matrix}.
Вводятся обозначения:
\begin{itemize}
    \item $X, Y$ -- матрицы размерностью $T \times N$ и $T \times M$ соответственно;
    \item $x_u, y_i$ -- столбцы $u, i$ матриц $X, Y$ соответственно, являющиеся представлением пользователя $u$ и объекта $i$;
    \item $R$ -- множество пар $(u, i)$ индексов пользователя и объекта, таких что оценка пользователя $u$ для объекта $i$ известна;
    \item $r$ -- матрица предпочтений, где $r_{u,i}$ оценка, которую пользователь $u$ выставил объекту $i$, $u \in \overline{1, N}, i \in \overline{1, M}$.
\end{itemize}
Оценка пользователем $u$ объекта $i$ вычисляется по формуле~\cite{yandex_matrix}:
\begin{equation}
    \hat{r}_{u,i} = x_u^Ty_i.
\end{equation}
Задача заключается в нахождении скрытых признаков.
Следовательно, необходимо минимизировать функцию потерь, в качестве которой выступает~\cite{yandex_matrix}:
\begin{equation}
    L(X, Y) = \sum_{(u, i) \in R}{(r_{ui} - \hat{r}_{u,i})^2}.
\end{equation}
Таким образом задача рекомендаций сведена к задачи оптимизации.
Тогда следует применить $L2$-регуляризацию, чтобы не произошло переобучение модели~\cite{yandex_matrix}.
После применения $L2$-регуляризации, функция потерь принимает вид:
\begin{equation}
    L(X, Y) = \sum_{(u, i) \in R}{(r_{ui} - \hat{r}_{u,i})^2} + \lambda \sum_{u = 1}^{N}{(||x_u||^2C_u)} + \lambda \sum_{i = 1}^{M}{(||y_i||^2C_i)},
    \label{eq:loss}
\end{equation} где $\lambda > 0$, а $C=(C_1,\dots,C_N)$ -- вектор констант.

Алгоритм минимизации функции потерь~\eqref{eq:loss} состоит из следующих шагов~\cite{yandex_matrix}:
\begin{itemize}
    \item начинается цикл вычислений до тех пор, пока функция потерь не примет достаточно маленькое значение;
    \item фиксируется матрица $X$;
    \item решается задача $L2$-регуляризованной регрессии для поиска матрицы $Y$;
    \item фиксируется матрица $Y$;
    \item решается задача $L2$-регуляризованной регрессии для поиска матрицы $X$;
    \item вычисляется новое значение функции потерь и выполняется переход в начало цикла.
\end{itemize}

Далее рассмотрен шаг алгоритма вычисления строки $x_u$ матрицы $X$.
Необходимо найти такое значение вектора $x_u$, для которого значение функции потерь~\eqref{eq:loss} минимально.
Тогда значение $x_u$ может быть представлено как~\cite{yandex_matrix}:
\begin{equation}
    \underset{x_u}{argmin} \sum_{(u, i) \in R}{(r_{ui} - x_u^Ty_i)^2} + \lambda \sum_{u = 1}^{N}{(||x_u||^2C_u)} + \lambda \sum_{i = 1}^{N}{(||y_i||^2C_i)},
    \label{eq:argmin_step_1}
\end{equation}
где функция $\underset{x_u}{argmin}$ обозначает такое $x_u$, при котором значение выражения минимально.
Выражение~\eqref{eq:argmin_step_1} упрощается~\cite{yandex_matrix}:
\begin{equation}
    \underset{x_u}{argmin} -2x_u^T\sum_{(u, i) \in R}{(r_{ui}y_i)} + x_u^T(\sum_{(u, i) \in R}{(y_iy_i^T) + \lambda C_uI})x_u,
    \label{eq:argmin_step_2}
\end{equation}
где $I$ -- единичная матрица.
Далее применяется следующее соотношение~\cite{yandex_matrix}:
\begin{equation}
    \underset{x_u}{argmin} -2x_u^TB + x_u^TAx_u = A^{-1}B.
    \label{eq:th1}
\end{equation}
Из~\eqref{eq:argmin_step_2}--\eqref{eq:th1} следует, что
\begin{equation}
    x_u = (\sum_{i: (u, i) \in R}{(y_iy_i^T)}+ \lambda C_uI)^{-1}\sum_{i: (u, i) \in R}{(r_{u,i}y_i)}.
\end{equation}

\subsection{Гибридные рекомендательные методы}
Существуют алгоритмы рекомендаций, которые предсказывают оценку объекта пользователем на основе как матрицы оценок, так и его параметров.
Вводятся следующие обозначения:
\begin{itemize}
    \item $U = (u_i | i \in \overline{1,k})$ -- кортеж рассматриваемых пользователей;
    \item $u' \notin U$ -- пользователь, для которого необходимо сделать рекомендацию;
    \item $p_i, i \in \overline{1, n}$ -- вектор из $m$ параметров объекта $i$, который оценил пользователь $u'$;
    \item $r_i, i \in \overline{1,n}$ -- вектор, где $r_{i,j}, j \in \overline{1,k}$ -- оценка объекта $i$ пользователем $u_j$;
    \item $y$ -- вектор, где $y_i, i \in \overline{1,n}$ -- оценка объекта $i$ пользователем $u'$.
\end{itemize}
Создаётся матрица $X$, строка $i$ которой является конкатенацией векторов $p_i$ и $r_i$.
Аналогично для объектов, не оценённых $u'$, строится матрица $X'$.
Необходимо на основании данных $X, y$ найти вектор $y'$, где $y'_i$ -- оценка для объекта, представляемого вектором $X'_i$.
Для этого разрабатывается модель $M$, в результате обучения которой на данных $(X, y)$ вектор $y'$ вычисляется по формуле $y' \approx M(X')$~\cite{segaran}.
В качестве такой модели выступает одна из следующих~\cite{segaran}:
\begin{itemize}
    \item решающее дерево;
    \item линейная регрессия;
    \item $K$ ближайших соседей;
    \item логистическая регрессия;
    \item нейросеть.
\end{itemize}

Более того, допустимо комбинировать перечисленные модели для получения составной~\cite{segaran}.
Пусть есть $N$ моделей $M_1,\dots,M_N$, обученные на данных $(X, y)$.
В результате их применения к $X$ получаются результаты $q_1,\dots,q_N$, которые являются вариантами вектора $y$.
Аналогично $M_i(X') = q'_i, i \in \overline{1,N}$.
Таким образом, появляется необходимость в разработке другой модели $M_0$, которая после обучения на данных $(q_1,\dots,q_N,y)$ используется для вычисления $y'$ по формуле $y' \approx M_0(q'_1,\dots,q'_N)$.
Алгоритм работы всей составной модели содержит следующие шаги~\cite{segaran}:
\begin{itemize}
    \item принять на вход матрицу параметров объектов $X'$;
    \item вычислить $q_i = M_i(X'), i \in \overline{1,N}$;
    \item вычислить $y' = M_0(q_1,\dots,q_N)$.
\end{itemize}
Таким образом, получается комбинация моделей для наиболее точного предсказания оценок объектов выбранным пользователем.

Далее рассмотрены особенности гибридного рекомендательного метода.
Рассмотренный метод сводит задачу рекомендаций к задаче машинного или глубокого обучения.
Это позволяет разработать несколько различных моделей и выбрать наиболее точную из них в результате тестирования.
Под точностью подразумевается значение функции потерь, которая вводится для оценки качества модели на тестовом наборе данных.
Также гибридный метод рекомендаций способен использовать наиболее полный набор данных, что повышает точность результатов.
Но при этом не предусматривается решение проблемы <<холодного старта>>, которая заключается в том, что если пользователь новый, что и данных о его оценках либо относительно мало, либо нет вообще.
Таким образом, модель для рекомендаций обучается на меньшем количестве данных, что сказывается на точности рекомендаций.
Другая сложность гибридного рекомендательного метода заключается в том, что вектора с оценками $r_i, i \in \overline{1,n}$ разрежены, так как каждый объект вероятнее всего не оценён каждым пользователем $i \in \overline {1,k}$.
Таким образом, данные, на которых обучается модель, не всегда достаточны для её корректной работы.
Следовательно, необходимо выставлять пропущенные оценки самостоятельно за других пользователей.
Пропущенная оценка может быть определена по умолчанию, случайным образом или с помощью какого-либо другого рекомендательного метода.